\documentclass{article}

\usepackage{makecell}

\title{Ch. 22 CW \#2}
\author{Principles of Microeconomics}
\date{\today}

\begin{document}

\maketitle

\subsection*{Instructions} Please justify all of your answers and submit your solutions on Gradescope by
the beginning of next class.

\subsection*{Problems}

\begin{enumerate}

\item For each of the following situations, draw a diagram containing three of Isabella's indifference curves. Be sure to label and number the axes of your diagrams.

	\begin{enumerate}
	
	\item For Isabella, cars and tires are perfect complements, but in a ratio of 1:4. That is, for each car, Isabella wants exactly four tires. Place tires on the horizontal axis and cars on the vertical axis.
	
	\item Isabella gets utility only from her caffeine intake. She can consume Valley Dew or cola, and Valley Dew contains twice as much caffeine as cola. Place cola on the horizontal axis and Valley Dew on the vertical axis.
	
	\item Isabella gets utility from two things: leisure time and consumption. Both have diminishing marginal utility. Place leisure on the horizontal axis and consumption on the vertical axis. You don't need to number your axes for this one.
	
	\item Isabella can consume two goods -- skis and bindings. For each ski, she wants exactly one binding. Place bindings on the horizontal axis and skis on the vertical axis.
	
	\item Isabella gets utility from consuming soda, but she gets no utility from consuming water. Water neither increases or decreases her utility. Place water on the horizontal axis and soda on the vertical axis.
	
	\end{enumerate}

\item Marta drinks three cups of coffee a day, for which she has diminishing marginal utility. which of her three cups generates the greatest increase in total utility? Which generates the least?
	
\item Suppose that Bridget and Erin spend their incomes on two goods, food ($F$) and clothing ($C$). Bridget's preferences are represented by the utility function $U(F, C) = 10FC$. Erin's preferences are represented by the utility function $U(F, C) = 0.2F^2C^2$. 

	\begin{enumerate}
	
	\item With food on the horizontal axis and clothing on the vertical axis, identify on a graph the set of points that give Bridget the same level of utility as the bundle $(10, 5)$. Do the same for Erin on a separate graph.
	
	\item On the same two graphs, identify the set of bundles that give Bridget and Erin the same level of utility as the bundle $(15, 8)$.
	
	\item Do you think Bridget and Erin have the same preferences or different preferences? Explain.
	
	\end{enumerate}
	
\item Julius has \$200 which he wants to spend on movies and ties. Movies cost \$20, and ties cost \$4.

	\begin{enumerate}
	
	\item Julius chooses to consume 4 movies and 30 ties. Draw this scenario with movies on the horizontal axis. What is Julius's marginal rate of substitution at his chosen basket?
	
	\item Julius's budget doubles to \$400. Draw this scenario. Explain what happens to Julius's quantity consumed of each good. State any assumptions that you're making.
	
	\item Julius's budget goes back to \$200, and the price of ties increases to \$6. Draw this scenario. How would his consumption bundle change relative to part (a)?
	
	\end{enumerate}
	
\item Bernie can spend his income on two different goods: notebooks and pens. The accompanying table shows Bernie's utilities from notebooks and pens.

	\begin{center}
	\begin{tabular}{c c | c c}
	\makecell{Quantity of\\Notebooks} & \makecell{Utility from\\Notebooks} & \makecell{Quantity of\\Pens} & \makecell{Utility from\\Pens} \\
	\hline
	0 & 0 & 0 & 0 \\
	1 & 32 & 2 & 28 \\
	2 & 60 & 4 & 52 \\
	3 & 84 & 6 & 72 \\
	4 & 104 & 8 & 96 \\
	5 & 120 & 10 & 108
	\end{tabular}
	\end{center}
	
	The price of a notebook is \$4, the price of a pen is \$2, and Bernie has \$20 of income to spend.
	
	\begin{enumerate}
	
	\item Which consumption bundles of notebooks and pens can Bernie consume if he spends all his income? Illustrate Bernie's budget line with a diagram, putting pens on the horizontal axis and notebooks on the vertical axis.
	
	\item Calculate the marginal utility of each notebook and the marginal utility of each pen. Then calculate the marginal utility per dollar spent on notebooks and the marginal utility per dollar spent on pens.
	
	\item Which bundle will Bernie choose?
	
	\end{enumerate}

\end{enumerate}

\end{document}